\phantomsection
\addcontentsline{toc}{subsection}{Unidad 1. Estequiometría en reacciones completas}
\subsection*{\textcolor{blueunam}{Unidad 1. Estequiometría en reacciones completas}}

\begin{notadidactica}
\textbf{Sobre las notas en este documento.} Hay tres niveles según su importancia:

\begin{itemize}
    \item \textcolor{rojoNota}{\textbf{Notas relevantes}} (línea roja gruesa): 
    convenciones, definiciones y sistemas de notación que se usan en todo el 
    documento. No se pueden ignorar.
    \item \textcolor{amarilloNota}{\textbf{Notas importantes}} (línea amarilla 
    media): aclaraciones conceptuales, justificaciones y puentes entre ideas. 
    Opcionales, pero recomendadas.
    \item \textcolor{azulNota}{\textbf{Notas didácticas}} (línea azul delgada): 
    analogías, ejemplos cotidianos e intuiciones. Totalmente opcionales.
\end{itemize}
\end{notadidactica}

\phantomsection
\addcontentsline{toc}{subsubsection}{Clase 1 -- 18/08/2026}
\subsubsection*{\textcolor{blueunam}{Clase 1 -- 18/08/2026: Definiciones y conceptos básicos}}

\begin{center}
\small
\begin{tabular}{|p{7cm}|p{7cm}|}
\hline
\textbf{Submicroscópico} & \textbf{Macroscópico} \\
\hline
Entidades químicas: objeto individual, partícula o conjunto de partículas enlazadas. &
Sustancias químicas: conjunto de entidades químicas iguales. \\[6pt]
Iones, átomos, moléculas & Sustancias químicas \\[6pt]
\ce{O2} oxígeno molecular & Tanque de oxígeno (\ce{O2}) \\
\hline
\end{tabular}
\normalsize
\end{center}

\textbf{Elemento:} Conjunto de núcleos idénticos.

\textbf{Mezcla:} Combinación de dos o más sustancias químicas.
\begin{itemize}
    \item Homogénea: una sola fase.
    \item Heterogénea: más de una fase.
\end{itemize}

\textbf{Disolución:} Mezcla homogénea de un soluto y un disolvente.

\textbf{Soluto:} Especie química que se estudia.

\textbf{Disolvente:} Especie química que disuelve al soluto.

\textbf{Cómo preparar una disolución:}
\begin{itemize}
    \item Soluto sólido (debe ser soluble): pesar $\to$ disolver en disolvente $\to$ completar hasta volumen deseado.
    \item Soluto líquido (debe ser miscible): medir volumen $\to$ disolver en disolvente $\to$ completar hasta volumen deseado.
\end{itemize}

\textbf{Cantidad de sustancia:}
\begin{equation}
    n = \frac{m}{MM}
    \quad \left[\frac{\mathrm{g}}{\mathrm{g/mol}} = \mathrm{mol}\right]
\end{equation}
donde $n$ es la cantidad de sustancia ($\mathrm{mol}$), $m$ la masa ($\mathrm{g}$) y $MM$ la masa molar ($\mathrm{g/mol}$).

\textbf{El mol como una abreviación}

Un mol no es más que una \textbf{abreviación}. Igual que abreviemos 
$1000\,\mathrm{g}$ a $1\,\mathrm{kg}$, o $1000\,\mathrm{mL}$ a $1\,\mathrm{L}$, 
el mol abrevia cantidad de sustancia. No abrevia masa ni volumen: abrevia 
número de objetos, cosas, partículas o moléculas.

Imagina un ejemplo cotidiano. Cuando pides un cartón de huevos, no pides 12 
huevos por separado: pides el cartón. O una caja que contiene 20 libros: no 
pides los 20 libros uno por uno, pides la caja. Lo mismo aplica al mol. La 
diferencia con estos ejemplos es que el mol no es una cantidad de 12 o de 20, 
sino de \num{6.022e23}. Pero que el número no te espante: es lo mismo, 
solo con una cifra muy grande.

Los químicos usan esto porque sería tedioso decir \num{6.022e23} 
moléculas de algo cada vez. En su lugar lo abrevian a $1\,\mathrm{mol}$. Si 
tuvieras una colección de tazos de esa cantidad, podrías decir que tienes 
$1\,\mathrm{mol}$ de tazos. No rompes con la definición: se puede hablar de lo 
que sea en mol, mientras se pueda contar.

La cuestión es que cotidianamente manejamos números mucho menores. Si tuvieras 
100 tazos, esa cantidad expresada en mol sería una fracción absurdamente 
pequeña, equivalente a $\num{1.6e-22}\,\mathrm{mol}$. Pero la idea se mantiene: el mol 
sirve para contar cualquier cosa, solo que en el día a día usamos cantidades 
tan pequeñas que la cifra en mol sale minúscula. Simplemente es una abreviación 
para no tener que decir números enormes al hablar de partículas o moléculas. 
Estas vienen en tantísimas partículas que ponerlas con números sin abreviación 
nos llevaría muchísimo más tiempo que simplemente abreviarlo.

\medskip
Volviendo con el ejemplo de las cajas de libros. Sabemos que cada caja tiene 
20 libros, pero no todos los libros son iguales. Pueden ser enciclopedias o 
cuentos infantiles. 20 enciclopedias pesan mucho más que 20 cuentos infantiles.

Lo importante es que todas las cajas tienen el \textbf{mismo número de libros} 
(20), sin importar qué libros contengan. Lo que cambia es cuánto pesa cada 
libro por separado.

Imaginemos que tenemos 5 cajas de enciclopedias. Si cada caja tiene una masa 
de $1000\,\mathrm{g}$, las 5 cajas juntas pesan $5000\,\mathrm{g}$.

\medskip
Ahora llevemos esto al mol. Es el mismo principio: el mol es la caja, y 
cambiamos libros por partículas.

El mol, como dijimos, solo es un número. $1\,\mathrm{mol}$ de algo puede pesar 
distinto según de qué se trate. $\text{mol} \neq \text{masa}$, y ahí está la 
clave: \textbf{la masa no es del mol, es de lo que el mol contiene}. El número 
20 es de cantidad, no de peso. Con el mol pasa igual: cuenta partículas, no 
mide masa.

Pero si sabemos cuánta masa tiene $1\,\mathrm{mol}$ de algo, podemos saber 
cuánta masa tendrá en $5\,\mathrm{mol}$ o en $0.5\,\mathrm{mol}$. Es una 
proporción directa: si $1\,\mathrm{mol}$ de algo tiene masa $M$, entonces 
$5\,\mathrm{mol}$ tienen masa $5M$, y $0.5\,\mathrm{mol}$ tienen masa $0.5M$.

\medskip
Una vez visto esto, podemos pensar en otra cosa. ¿Por qué el mol nos hace 
pasar de una unidad de partículas a masa?

La respuesta es más simple de lo que parece. Si siempre tomas el mismo número 
de partículas, la masa que obtienes depende de cuánto pesa cada una. Un átomo 
de hierro pesa más que uno de carbono. Entonces $1\,\mathrm{mol}$ de hierro 
pesa más que $1\,\mathrm{mol}$ de carbono, aunque ambos tengan el mismo número 
de átomos. \textbf{El mol es solo un paquete de partículas; la masa del paquete 
depende del contenido.}

\medskip
Con eso en mente, veamos ejemplos concretos.

Imagina $1\,\mathrm{mol}$ de carbono puro. Sabemos que son muchísimas 
partículas, pero ¿cuánto pesa eso? Casi $12\,\mathrm{g}$. Y $12\,\mathrm{g}$ de 
cualquier cosa son perceptibles a simple vista: no importa qué forma te 
imagines, siempre los verás. Lo mismo con $1\,\mathrm{mol}$ de hierro: casi 
$56\,\mathrm{g}$. Una barra de metal de ese peso la ves sin problema.

No todos los elementos funcionan igual. $1\,\mathrm{mol}$ de hidrógeno pesa 
solo $2\,\mathrm{g}$, y $2\,\mathrm{g}$ de gas no los ves. El carbono y el 
hierro sí funcionan como ejemplo porque sus masas molares dan cantidades 
perceptibles.

La cuestión es que ambos son apreciables a simple vista, y lo curioso es que 
ambos están compuestos por el mismo número de partículas. Es decir, contienen 
el mismo número de átomos, pero de elementos distintos.

\medskip
Aquí es donde uno podría tener el concepto erróneo de que un mol es 
inimaginable. \num{6.022e23} partículas son muchísimas, sí, pero son 
pequeñas. Es como imaginar 5000 hormigas. Son muchísimas, pero cada una mide 
1 cm. Si las pones en fila, ocupan 50 metros, algo que ves en tu campo de 
visión sin problema. Si en vez de 5000 hormigas fueran 5000 personas, la cosa 
cambia: una fila de 5000 personas ocuparía más de 8 kilómetros. \textit{Mismo 
número, percepción completamente distinta, porque el tamaño de cada unidad 
importa tanto como la cantidad.}

Con las partículas pasa algo similar. ¿Son tantísimas como para que un 
compuesto formado por ese número sea apreciable? La respuesta la tenemos aquí: 
sí, es suficiente para ser apreciado a simple vista, pero no tantísimo como 
para pensar que $1\,\mathrm{mol}$ de carbono llenaría tu casa.

\medskip
Aquí hay que aclarar algo: depende del compuesto. El número de partículas de 
$1\,\mathrm{mol}$ es el mismo sin importar de qué hablemos. La diferencia está 
en la masa y en el volumen que ocupa. Decíamos antes que $1\,\mathrm{mol}$ de 
hierro tiene más masa que $1\,\mathrm{mol}$ de carbono ($56\,\mathrm{g}$ vs 
$12\,\mathrm{g}$). ¿Y el volumen? $56\,\mathrm{g}$ de hierro ocupan 
aproximadamente $7\,\mathrm{mL}$, mientras que $12\,\mathrm{g}$ de carbono 
(grafito) ocupan aproximadamente $5\,\mathrm{mL}$. El hierro ocupa un poco 
más, aunque la diferencia no es enorme. \textit{Masa y volumen son propiedades 
distintas, y la densidad las conecta.}

En cualquier caso, es bueno tener esa noción y estar familiarizados con el 
mol, no verlo como algo extraño. En resumen, el mol es una cantidad de algo, 
igual que una docena o un cartón de huevos. Fácilmente podría ser como 
contamos nuestra colección de tazos. No es más que eso: una abreviación 
para no escribir números enormes. \\

\begin{notafija}
\textbf{Convención de subíndices:} en las expresiones matemáticas se usan las 
siguientes abreviaturas en subíndices: $\mathrm{sol}$ = soluto, 
$\mathrm{dis}$ = disolvente, $\mathrm{dil}$ = disolución.
\end{notafija}

\textbf{Concentración en peso:}
\begin{align*}
    \rho_{\mathrm{sol}} &= \frac{m_{\mathrm{sol}}}{V_{\mathrm{dil}}}
    \quad \left[\frac{\mathrm{g}}{\mathrm{mL}}\right]
\end{align*}
En general es una relación \(m/V\); las unidades más comunes son $\mathrm{g/mL}$, $\mathrm{g/L}$, $\mathrm{mg/mL}$ y $\mathrm{mg/L}$. Lo más frecuente es $\mathrm{g/mL}$.

\begin{notafija}
La notación \(c\) todavía se usa para concentración en peso, pero por su 
ambigüedad con la molaridad se omite y se usa \(\rho\) (recomendación IUPAC 
para concentración de masa).
\end{notafija}

\textbf{Concentración porcentual:}
\begin{align*}
    \%\,m/m &= \frac{m_{\mathrm{sol}}[\mathrm{g}]}{m_{\mathrm{dil}}[\mathrm{g}]} \times 100 \quad [\%] \\[4pt]
    \%\,m/v &= \frac{m_{\mathrm{sol}}[\mathrm{g}]}{V_{\mathrm{dil}}[\mathrm{mL}]} \times 100 \quad \left[\frac{\mathrm{g}}{100\,\mathrm{mL}}\right] \\[4pt]
    \%\,v/v &= \frac{V_{\mathrm{sol}}[\mathrm{mL}]}{V_{\mathrm{dil}}[\mathrm{mL}]} \times 100 \quad [\%]
\end{align*}

\textbf{Partes por millón (ppm):}
\begin{align*}
    \mathrm{ppm}\,(m/m) &\equiv \frac{m_{\mathrm{sol}}[\mathrm{g}]}{m_{\mathrm{dil}}[\mathrm{g}]} \times 10^6 \quad \left[\frac{\mathrm{mg}}{\mathrm{kg}}\right] \\[4pt]
    \mathrm{ppm}\,(m/v) &\equiv \frac{m_{\mathrm{sol}}[\mathrm{g}]}{V_{\mathrm{dil}}[\mathrm{mL}]} \times 10^6 \quad \left[\frac{\mathrm{mg}}{\mathrm{L}}\right] \\[4pt]
    \mathrm{ppm}\,(v/v) &\equiv \frac{V_{\mathrm{sol}}[\mathrm{mL}]}{V_{\mathrm{dil}}[\mathrm{mL}]} \times 10^6 \quad \left[\frac{\mu\mathrm{L}}{\mathrm{L}}\right]
\end{align*}

\textbf{Masa de la disolución:}
\begin{align*}
    m_{\mathrm{dil}}[\mathrm{g}] &= m_{\mathrm{sol}}[\mathrm{g}] + m_{\mathrm{dis}}[\mathrm{g}]
\end{align*}

La masa siempre es aditiva. El volumen no siempre lo es:
\begin{itemize}
    \item Caso 1: misma naturaleza soluto/disolvente, concentraciones similares: volumen aditivo.
    \item Caso 2: dos sustancias no miscibles: volumen aditivo.
\end{itemize}

\phantomsection
\addcontentsline{toc}{subsubsection}{Ejercicios de concentración}
\subsubsection*{\textcolor{blueunam}{Ejercicios de concentración}}

\begin{exercise}{1}{}
Calcula el $\%\,m/m$ si se tiene una disolución de $50.0\,\mathrm{mL}$ cuya densidad es $0.762\,\mathrm{g/mL}$ y se sabe que hay disuelto $0.0500\,\mathrm{mol}$ de \ce{CoSO4}.

Solución:
\begin{align*}
    &\%\,m/m = ? \\[4pt]
    &V_{\mathrm{dil}} = 50.0\,\mathrm{mL} \\[4pt]
    &\rho_{\mathrm{dil}} = 0.762\,\mathrm{g/mL} \\[4pt]
    &n_{\mathrm{sol}} = 0.0500\,\mathrm{mol}\;\ce{CoSO4} \\[4pt]
    &MM_{\mathrm{CoSO_4}} = 58.93 + 32.07 + 4(16.00) = 154.99\,\mathrm{g/mol} \\[4pt]
    &m_{\mathrm{dil}} = V_{\mathrm{dil}} \cdot \rho_{\mathrm{dil}}
        = (50.0\,\mathrm{mL})(0.762\,\mathrm{g/mL})
        = 38.1\,\mathrm{g}_{\mathrm{dil}} \\[4pt]
    &m_{\mathrm{sol}} = n_{\mathrm{sol}} \cdot MM_{\mathrm{CoSO_4}}
        = (0.0500\,\mathrm{mol})(154.99\,\mathrm{g/mol})
        = 7.750\,\mathrm{g}_{\mathrm{sol}} \\[4pt]
    &\%\,m/m = \frac{m_{\mathrm{sol}}}{m_{\mathrm{dil}}} \times 100
        = \frac{7.750\,\mathrm{g}_{\mathrm{sol}}}{38.1\,\mathrm{g}_{\mathrm{dil}}} \times 100
        = 20.34\%\,m/m
\end{align*}
\end{exercise}

\begin{exercise}{2}{}
Calcula la concentración en partes por millón si se disuelven $\num{5.000e-3}\,\mathrm{g}$ de bromuro de calcio en $200.0\,\mathrm{mL}$ de agua. Considera la densidad del agua como $1\,\mathrm{g/mL}$.

Solución:
\begin{align*}
    &\mathrm{ppm} = ? \\[4pt]
    &m_{\mathrm{sol}} = \num{5.000e-3}\,\mathrm{g}\;\ce{CaBr2} \\[4pt]
    &V_{\mathrm{dis}} = 200.0\,\mathrm{mL} \\[4pt]
    &\rho_{\mathrm{dis}} = 1\,\mathrm{g/mL} \\[4pt]
    &m_{\mathrm{dis}} = V_{\mathrm{dis}} \cdot \rho_{\mathrm{dis}}
        = (200.0\,\mathrm{mL})(1\,\mathrm{g/mL})
        = 200.0\,\mathrm{g}_{\mathrm{dis}} \\[4pt]
    &m_{\mathrm{dil}} = m_{\mathrm{sol}} + m_{\mathrm{dis}}
        = \num{5.000e-3}\,\mathrm{g} + 200.0\,\mathrm{g}
        = 200.005\,\mathrm{g}_{\mathrm{dil}} \\[4pt]
    &\mathrm{ppm} = \frac{m_{\mathrm{sol}}}{m_{\mathrm{dil}}} \times 10^6
        = \frac{\num{5.000e-3}\,\mathrm{g}}{200.005\,\mathrm{g}} \times 10^6
        = 25.00\,\mathrm{ppm}
\end{align*}
\end{exercise}

\begin{exercise}{3}{}
¿Qué concentración de glucosa posee una disolución inicial si para preparar $125\,\mathrm{mL}$ de disolución acuosa de glucosa al $3\%\,m/v$ se han añadido $50\,\mathrm{mL}$ de agua? Considera que los volúmenes se adicionan.

Solución:
\begin{align*}
    &\%\,m/v_{\mathrm{i}} = ? \\[4pt]
    &V_{\mathrm{f}} = 125\,\mathrm{mL} \\[4pt]
    &\%\,m/v_{\mathrm{f}} = 3\%\,m/v \\[4pt]
    &V_{\mathrm{+agua}} = 50\,\mathrm{mL}
\end{align*}
Análisis: la disolución \textit{final} de $125\,\mathrm{mL}$ al $3\%\,m/v$ contiene $3\,\mathrm{g}$ de soluto por cada $100\,\mathrm{mL}$. La masa de soluto es la misma antes y después de agregar agua (solo se diluyó). El volumen inicial es el final menos el agua añadida.
\begin{align*}
    &m_{\mathrm{sol}} = \frac{3\,\mathrm{g}_{\mathrm{sol}}}{100\,\mathrm{mL}_{\mathrm{dil}}} \times 125\,\mathrm{mL}_{\mathrm{dil}} = 3.75\,\mathrm{g}_{\mathrm{sol}} \\[4pt]
    &V_{\mathrm{i}} = V_{\mathrm{f}} - V_{\mathrm{+agua}} = 125\,\mathrm{mL} - 50\,\mathrm{mL} = 75\,\mathrm{mL} \\[4pt]
    &\%\,m/v_{\mathrm{i}} = \frac{m_{\mathrm{sol}}}{V_{\mathrm{i}}} \times 100
        = \frac{3.75\,\mathrm{g}_{\mathrm{sol}}}{75\,\mathrm{mL}} \times 100
        = 5.00\%\,m/v
\end{align*}
\end{exercise}

\begin{exercise}{4}{}
Calcula la concentración en peso $(\mathrm{g/mL})$ y el $\%\,m/v$ de una disolución acuosa de $100.0\,\mathrm{mL}$ que contiene $0.025\,\mathrm{mol}$ de \ce{H2CO3}.

Solución:
\begin{align*}
    &\mathrm{g/mL} = ? \\[4pt]
    &\%\,m/v = ? \\[4pt]
    &V_{\mathrm{dil}} = 100.0\,\mathrm{mL} \\[4pt]
    &n_{\mathrm{sol}} = 0.025\,\mathrm{mol}\;\ce{H2CO3} \\[4pt]
    &MM_{\mathrm{H_2CO_3}} = 62.03\,\mathrm{g/mol} \\[4pt]
    &m_{\mathrm{sol}} = MM_{\mathrm{H_2CO_3}} \cdot n_{\mathrm{sol}}
        = (62.03\,\mathrm{g/mol})(0.025\,\mathrm{mol}\;\ce{H2CO3})
        = 1.551\,\mathrm{g}_{\mathrm{sol}} \\[4pt]
    &g/mL = \frac{m_{\mathrm{sol}}}{V_{\mathrm{dil}}}
        = \frac{1.551\,\mathrm{g}_{\mathrm{sol}}}{100.0\,\mathrm{mL}}
        = 0.01551\,\mathrm{g/mL} \\[4pt]
    &\%\,m/v = g/mL \times 100
        = (0.01551\,\mathrm{g/mL}) \times 100
        = 1.551\%\,m/v
\end{align*}
\end{exercise}

\phantomsection
\addcontentsline{toc}{subsubsection}{Clase 2 -- 25/08/2026}
\subsubsection*{\textcolor{blueunam}{Clase 2 -- 25/08/2026: Expresiones de la concentración en unidades Químicas}}

\textbf{Molaridad:}
\begin{align*}
    c_{\mathrm{sol}} &= \frac{n_{\mathrm{sol}}[\mathrm{mol}]}{V_{\mathrm{dil}}[\mathrm{L}]}
    \quad \left[\frac{\mathrm{mol}}{\mathrm{L}}\right]
\end{align*}

\begin{notafija}
La notación \(M\) todavía se usa, pero por su ambigüedad con la masa molar se 
omite en estos apuntes.
\end{notafija}

\textbf{Molalidad:}
\begin{align*}
    b_{\mathrm{sol}} &= \frac{n_{\mathrm{sol}}[\mathrm{mol}]}{m_{\mathrm{dis}}[\mathrm{kg}]}
    \quad \left[\frac{\mathrm{mol}}{\mathrm{kg}}\right]
\end{align*}

\begin{notafija}
La notación \(m\) todavía se usa, pero por su ambigüedad con la masa se omite 
en estos apuntes.
\end{notafija}

\textbf{Equivalente:} cantidad de sustancia que reacciona con, o aporta, 
exactamente 1 mol de una unidad reactiva (H$^{\color{catred}+}$, 
OH$^{\color{anblue}-}$, e$^{\color{anblue}-}$ o carga, según el tipo de reacción).

\textbf{Factor de equivalencia $z$:} número de equivalentes por mol de especie.
\begin{align*}
    z &= \frac{n_{eq}}{n_{\mathrm{sol}}}
    \quad \left[\frac{\mathrm{eq}}{\mathrm{mol}}\right] \\[4pt]
    n_{eq} &= z \cdot n_{\mathrm{sol}}
    \quad [\mathrm{eq}]
\end{align*}

donde:
\begin{itemize}
    \item \(n_{eq}\) es la cantidad de equivalentes, en eq
    \item \(n_{\mathrm{sol}}\) es la cantidad de sustancia, en mol
    \item \(z\) es el factor de equivalencia, en eq/mol
\end{itemize}

\textbf{Equivalente ácido:} cantidad de ácido que aporta 1 mol de H$^{\color{catred}+}$.
\begin{align*}
    \ce{HCl} &\longrightarrow \ce{H}^{\color{catred}+} + \ce{Cl}^{\color{anblue}-} 
    & z &= 1\,\mathrm{eq/mol} \\[4pt]
    \ce{H2SO4} &\longrightarrow 2\,\ce{H}^{\color{catred}+} + \ce{SO4}^{2\color{anblue}-} 
    & z &= 2\,\mathrm{eq/mol} \\[4pt]
    \ce{H3PO4} &\longrightarrow 3\,\ce{H}^{\color{catred}+} + \ce{PO4}^{3\color{anblue}-} 
    & z &= 3\,\mathrm{eq/mol}
\end{align*}

\textbf{Equivalente básico:} cantidad de base que aporta 1 mol de 
OH$^{\color{anblue}-}$.
\begin{align*}
    \ce{NaOH} &\longrightarrow \ce{Na}^{\color{catred}+} + \ce{OH}^{\color{anblue}-} 
    & z &= 1\,\mathrm{eq/mol} \\[4pt]
    \ce{Ca(OH)2} &\longrightarrow \ce{Ca}^{2\color{catred}+} + 2\,\ce{OH}^{\color{anblue}-} 
    & z &= 2\,\mathrm{eq/mol} \\[4pt]
    \ce{Al(OH)3} &\longrightarrow \ce{Al}^{3\color{catred}+} + 3\,\ce{OH}^{\color{anblue}-} 
    & z &= 3\,\mathrm{eq/mol}
\end{align*}

\textbf{Equivalente redox:} cantidad de especie que transfiere 1 mol de 
e$^{\color{anblue}-}$.
\begin{align*}
    \ce{Na} &\longrightarrow \ce{Na}^{\color{catred}+} + \ce{e}^{\color{anblue}-} 
    & z &= 1\,\mathrm{eq/mol} \\[4pt]
    \ce{Mg} &\longrightarrow \ce{Mg}^{2\color{catred}+} + 2\,\ce{e}^{\color{anblue}-} 
    & z &= 2\,\mathrm{eq/mol} \\[4pt]
    \ce{Al} &\longrightarrow \ce{Al}^{3\color{catred}+} + 3\,\ce{e}^{\color{anblue}-} 
    & z &= 3\,\mathrm{eq/mol}
\end{align*}

\textbf{Normalidad:}
\begin{align*}
    N &= \frac{n_{eq}[\mathrm{eq}]}{V_{\mathrm{dil}}[\mathrm{L}]}
    = z \cdot c_{\mathrm{sol}}
    \quad \left[\frac{\mathrm{eq}}{\mathrm{L}}\right]
\end{align*}

\begin{notaimportante}
\(n_{eq}\) es la cantidad de equivalentes [eq], \(n_{\mathrm{sol}}\) es la 
cantidad de sustancia [mol], y \(z\) es el factor de equivalencia [eq/mol]. 
La relación entre ellos es:
\[
n_{eq} = z \cdot n_{\mathrm{sol}}
\]
La IUPAC deprecó la normalidad y recomienda \(c_{\mathrm{eq}} = z \cdot 
c_{\mathrm{sol}}\) con \(z\) en [eq/mol]. Sin embargo, \(N\) sigue siendo común 
en cursos, valoraciones y reactivos comerciales, por lo que se usa aquí.
\end{notaimportante}

\textbf{Dependencia de la reacción:} \(z\) no es propiedad fija de una sustancia. 
Depende de la reacción específica en la que participa. Por eso una misma disolución 
puede tener distinta normalidad según el medio. Ejemplo con \ce{KMnO4}:

\begin{align*}
    \text{Medio ácido:}\quad \ce{MnO4}^{\color{anblue}-} + 8\,\ce{H}^{\color{catred}+} + 5\,\ce{e}^{\color{anblue}-} 
    &\longrightarrow \ce{Mn}^{2\color{catred}+} + 4\,\ce{H2O} & z &= 5\,\mathrm{eq/mol} \\[4pt]
    \text{Medio neutro:}\quad \ce{MnO4}^{\color{anblue}-} + 2\,\ce{H2O} + 3\,\ce{e}^{\color{anblue}-} 
    &\longrightarrow \ce{MnO2}(s) + 4\,\ce{OH}^{\color{anblue}-} & z &= 3\,\mathrm{eq/mol} \\[4pt]
    \text{Medio básico:}\quad \ce{MnO4}^{\color{anblue}-} + \ce{e}^{\color{anblue}-} 
    &\longrightarrow \ce{MnO4}^{2\color{anblue}-} & z &= 1\,\mathrm{eq/mol}
\end{align*}

Por eso al reportar la normalidad hay que especificar \textbf{para qué reacción 
fue preparada}. Decir ``disolución 2 N de \ce{KMnO4}'' sin especificar el medio es 
ambiguo. Lo correcto: ``disolución de \ce{KMnO4} 2 N en medio ácido'' (o el medio 
que corresponda).

\begin{notaimportante}
La IUPAC no reconoce el equivalente como unidad del SI; estrictamente, \(z\) es 
adimensional y \(n_{eq}\) se mide en mol. Sin embargo, tratar eq como unidad 
permite distinguir la normalidad (eq/L) de la molaridad (mol/L), y es la 
convención usada en el curso. Por eso se mantiene eq y eq/mol en estos apuntes.
\end{notaimportante}

\medskip
\textbf{Fracción molar:}
\begin{align*}
    \chi_i &= \frac{n_i[\mathrm{mol}]}{n_{\mathrm{total}}[\mathrm{mol}]}
    \quad [\text{adimensional}] \\[4pt]
    n_{\mathrm{total}} &= \sum_{i=1}^{k} n_i
    \quad \text{($k$ = número de componentes)} \\[4pt]
    \sum_{i=1}^{k} \chi_i &= 1
\end{align*}

\begin{notafija}
\textbf{Propiedades extensivas e intensivas.}

En química hay dos tipos de propiedades: las que dependen de cuánta sustancia 
tengas y las que no.

\textbf{Propiedades extensivas:} dependen de la cantidad. Si duplicas la 
cantidad de sustancia, la propiedad se duplica. Ejemplos: masa, volumen, mol, 
número de partículas.

\textbf{Propiedades intensivas:} no dependen de la cantidad. Si duplicas la 
cantidad de sustancia, la propiedad no cambia. Ejemplos: densidad, temperatura, 
concentración, punto de fusión.

\medskip
Imagina un vaso de agua y una cubeta de agua, ambas a la misma temperatura. La 
masa y el volumen de la cubeta son mucho mayores que los del vaso (extensivas). 
Pero la temperatura es la misma (intensiva), y la densidad también (intensiva). 
No importa cuánta agua tomes: su temperatura y su densidad no cambian.

Lo mismo pasa con una disolución. La masa total y los mol totales cambian 
según cuánta tengas (extensivas). Pero la concentración es la misma en cualquier 
porción (intensiva). Si tomas un vaso de agua de limón, sabe igual que la jarra 
completa, porque la concentración es intensiva.

\medskip
Esta distinción es clave para entender por qué en algunos ejercicios podemos 
asumir un volumen arbitrario. Las propiedades intensivas no dependen del tamaño 
de la muestra, así que se pueden calcular para cualquier cantidad. Las 
extensivas no: si no sabes cuánta sustancia hay, no puedes calcular su masa 
total ni sus mol totales.
\end{notafija}

\phantomsection
\addcontentsline{toc}{subsubsection}{Clase 3 -- 27/08/2026}
\subsubsection*{\textcolor{blueunam}{Clase 3 -- 27/08/2026: Más ejercicios sobre expresiones de la concentración}}

\textbf{Alicuota {\color{red}{(alc)}}:} porción de una disolución que conserva las mismas 
propiedades intensivas que la disolución original. Específicamente: concentración 
y composición. Lo que cambia es la cantidad (volumen, masa, mol), no la 
concentración.

\begin{exercise}{1}{}
Se disuelven $0.61\,\mathrm{g}$ de sulfato de cobre pentahidratado, 
\ce{CuSO4.5H2O}, en agua hasta completar $200.0\,\mathrm{cm^3}$ de disolución. 
Calcula la molaridad de la disolución. Después, se toman $25.0\,\mathrm{cm^3}$ 
de la disolución anterior y se le añade más agua hasta completar 
$400.0\,\mathrm{cm^3}$. ¿Cuál será la molaridad de la nueva disolución? 
Describe, mediante un esquema, el procedimiento experimental.

Solución:
\begin{align*}
    &c_{\mathrm{sol_2}} = ? \\[4pt]
    &m_{\mathrm{sol_1}} = 0.61\,\mathrm{g}\;\ce{CuSO4.5H2O} \\[4pt]
    &V_{\mathrm{dil_1}} = 200.0\,\mathrm{cm^3} = 0.200\,\mathrm{L} \\[4pt]
    &V_{\mathrm{alc}} = 25.0\,\mathrm{cm^3} = 0.025\,\mathrm{L} \\[4pt]
    &V_{\mathrm{dil_2}} = 400.0\,\mathrm{cm^3} = 0.400\,\mathrm{L} \\[4pt]
    &MM_{\mathrm{\ce{CuSO4.5H2O}}} = 249.69\,\mathrm{g/mol} \\[4pt]
    &n_{\mathrm{sol_1}} = \frac{m_{\mathrm{sol_1}}}{MM_{\mathrm{\ce{CuSO4.5H2O}}}}
    = \frac{0.61\,\mathrm{g}\;\ce{CuSO4.5H2O}}{249.69\,\mathrm{g/mol}}
    = \num{2.44e-3}\,\mathrm{mol}\;\ce{CuSO4.5H2O} \\[4pt]
    &c_{\mathrm{sol_1}} = c_{\mathrm{alc}} = \frac{n_{\mathrm{sol_1}}}{V_{\mathrm{dil_1}}}
    = \frac{\num{2.44e-3}\,\mathrm{mol}\;\ce{CuSO4.5H2O}}{0.200\,\mathrm{L}}
    = \num{1.22e-2}\,\mathrm{mol/L}
\end{align*}

\textbf{Análisis:}

El problema se resuelve en tres pasos, sin necesidad de fórmulas especiales:

\begin{enumerate}
    \item \textbf{mol en la alícuota.} La disolución 1 contiene 
    $n_{\mathrm{dil_1}} = \num{2.44e-3}\,\mathrm{mol}$ de \ce{CuSO4.5H2O} 
    en un volumen total $V_{\mathrm{dil_1}} = 200.0\,\mathrm{mL}$. Se toma una 
    alícuota de $V_{\mathrm{alc}} = 25.0\,\mathrm{mL}$. La cantidad de soluto en 
    la alícuota es proporcional al volumen tomado:
    \begin{align*}
        n_{\mathrm{alc}} &= n_{\mathrm{dil_1}} \cdot \frac{V_{\mathrm{alc}}}{V_{\mathrm{dil_1}}} \\[4pt]
        n_{\mathrm{alc}} &= (\num{2.44e-3}\,\mathrm{mol}) \cdot 
        \frac{25.0\,\mathrm{mL}}{200.0\,\mathrm{mL}} \\[4pt]
        n_{\mathrm{alc}} &= \num{3.05e-4}\,\mathrm{mol}
    \end{align*}
    Análisis dimensional: $\mathrm{mol} \cdot (\mathrm{mL}/\mathrm{mL}) = \mathrm{mol}$. 
    La razón de volúmenes es adimensional porque ambos están en la misma unidad.

    \item \textbf{Los mol no cambian al diluir.} Al agregar agua no se agrega ni 
    se quita soluto. Los $\num{3.05e-4}\,\mathrm{mol}$ de \ce{CuSO4.5H2O} 
    que había en la alícuota siguen siendo los mismos en la disolución 2.
    \begin{align*}
        n_{\mathrm{sol_2}} = n_{\mathrm{alc}} = \num{3.05e-4}\,\mathrm{mol}
    \end{align*}

    \item \textbf{Nueva concentración.} La misma cantidad de soluto ahora está 
    disuelta en un volumen mayor ($V_{\mathrm{dil_2}} = 400.0\,\mathrm{mL} = 0.400\,\mathrm{L}$). 
    Por definición de concentración:
    \begin{align*}
        c_{\mathrm{sol_2}} = \frac{n_{\mathrm{sol_2}}}{V_{\mathrm{dil_2}}}
        = \frac{\num{3.05e-4}\,\mathrm{mol}}{0.400\,\mathrm{L}}
        = \num{7.63e-4}\,\mathrm{mol/L}
    \end{align*}
\end{enumerate}

\textbf{Nota:} el cálculo también se puede resolver directamente con la relación 
$c_1 V_1 = c_2 V_2$, que es la versión comprimida de los tres pasos anteriores. 
En este problema, la alícuota actúa como ``disolución 1'' y la disolución final 
como ``disolución 2'':

\begin{align*}
    c_{\mathrm{alc}} V_{\mathrm{alc}} &= c_{\mathrm{sol_2}} V_{\mathrm{dil_2}} \\[4pt]
    c_{\mathrm{sol_2}} &= \frac{c_{\mathrm{alc}} \cdot V_{\mathrm{alc}}}{V_{\mathrm{dil_2}}}
    = \frac{(\num{1.22e-2}\,\mathrm{mol/L})(0.025\,\mathrm{L})}{0.400\,\mathrm{L}}
    = \num{7.63e-4}\,\mathrm{mol/L}
\end{align*}

La relación $c_1 V_1 = c_2 V_2$ funciona porque $n = c \cdot V$ y los mol 
de soluto son los mismos antes y después de diluir. Pero conviene recordar qué 
representa cada término:

\begin{itemize}
    \item La cantidad que se conserva es la de \textbf{soluto}, no la de 
    disolución. Al agregar agua, la cantidad de disolvente aumenta, el volumen 
    total aumenta, pero los mol de soluto no cambian.
    \item Por eso la concentración disminuye: al repartir la misma cantidad 
    de soluto en un volumen mayor, cada mililitro de disolución contiene menos 
    soluto.
    \item La masa de soluto también se conserva. Si se trabajara con \%-m/m, 
    la masa total sí cambiaría (porque se agrega agua), pero la masa de soluto 
    no. Es la misma idea en otra unidad.
\end{itemize}

\textit{En resumen:} al diluir, el soluto se redistribuye en un volumen mayor. 
La cantidad de soluto (mol o masa) es la misma; lo que cambia es el volumen 
de la disolución que lo contiene, y por lo tanto su concentración.

\textbf{Procedimiento experimental:}
\begin{enumerate}
    \item Pesar $0.61\,\mathrm{g}$ de \ce{CuSO4.5H2O} en balanza analítica.
    \item Transferir a matraz aforado de $200.0\,\mathrm{cm^3}$.
    \item Disolver con agua destilada y agitar.
    \item Aforar hasta la marca.
    \item Tomar alícuota de $25.0\,\mathrm{cm^3}$ con pipeta volumétrica.
    \item Transferir a matraz aforado de $400.0\,\mathrm{cm^3}$.
    \item Aforar hasta la marca y agitar.
\end{enumerate}
\end{exercise}

\begin{notadidactica}
\textbf{Nota conceptual: alícuota y dilución (ejemplo del pozole).}

\textit{Ejemplo didáctico.} Imaginemos un pozole bien revuelto en su olla. 
El caldo sería el disolvente y la carne el soluto. Aclaración: en una disolución 
verdadera el soluto está molecularmente disperso, no en trozos. Aquí la carne 
es una suspensión, pero para visualizar la idea de alícuota funciona.

Nos servimos una cuchara de pozole (alícuota) y la pasamos a nuestro plato. 
¿Cambió la concentración al pasar de la olla a la cuchara? No. ¿Cambió el sabor 
al comer directo de la olla vs. de la cuchara? No. La concentración es la misma 
en cualquier punto porque la olla está bien revuelta.

Si en la olla hay 40 piezas de pollo repartidas en 20 L de pozole, la 
concentración es:

\[
c = \frac{40\,\text{piezas}}{20\,\text{L}} = 2\,\text{piezas/L}
\]

Si la cuchara puede tomar 1 L de pozole, ¿cuántas piezas habrá en ella?

\[
n_{\text{cuchara}} = c \cdot V_{\text{cuchara}} = (2\,\text{piezas/L})(1\,\text{L}) = 2\,\text{piezas}
\]
\[
c_{\text{cuchara}} = \frac{n_{\text{cuchara}}}{V_{\text{cuchara}}} = \frac{2\,\text{piezas}}{1\,\text{L}} = 2\,\text{piezas/L}
\]

Aquí se ve que la concentración se conserva (2 piezas/L en la olla, 2 piezas/L 
en la cuchara), pero la cantidad de pollo no es la misma (40 piezas en la olla, 
2 piezas en la cuchara). Lo que se conserva es la fracción:

\[
\frac{40}{20} = \frac{2}{1} = 2
\]

Misma razón, números distintos.

Ahora pasamos la cuchara al plato y notamos que le falta caldo. Agregamos solo 
caldo, sin carne. ¿Cambió la concentración? Sí, pero no porque haya cambiado 
la cantidad de pollo. La cantidad de pollo sigue siendo 2 piezas. Lo que cambió 
es el volumen: antes era 1 L, ahora es 1 L + el caldo agregado. La nueva 
concentración es:

\[
c_{\text{nueva}} = \frac{n_{\text{pollo}}}{V_{\text{nuevo}}}
= \frac{2\,\text{piezas}}{V_{\text{plato}}}
\]

De aquí salen dos conclusiones:
\begin{itemize}
    \item La concentración baja si disminuye la cantidad de soluto.
    \item La concentración baja (más común) si aumenta el volumen del disolvente 
    sin agregar soluto.
\end{itemize}

En una disolución verdadera el soluto está molecularmente disperso, por eso la 
alícuota siempre es representativa. En una suspensión como el pozole, con trozos 
grandes, la cuchara podría no serlo por azar. Pero la idea de concentración 
intensiva y cantidad extensiva se mantiene.

Esta misma lógica se aplica en el ejercicio 6, donde se usa una \textbf{alícuota 
imaginaria} para simplificar los cálculos sin perder generalidad.
\end{notadidactica}

\begin{exercise}{6}{}
Calcula la molalidad de una disolución de THF (\ce{C4H8O}) 
$2.25\,\mathrm{mol/L}$ cuya densidad es $0.89\,\mathrm{g/mL}$. Considera la 
densidad del agua como $1\,\mathrm{g/mL}$.

Solución:

\begin{align*}
    &b_{\mathrm{sol}} = ? \\[4pt]
    &c_{\mathrm{sol}} = 2.25\,\mathrm{mol/L} \\[4pt]
    &\rho_{\mathrm{dil}} = 0.89\,\mathrm{g/mL} \\[4pt]
    &\rho_{\ce{H2O}} = 1\,\mathrm{g/mL}
\end{align*}

\begin{notaimportante}
\textbf{Sobre la base de cálculo (alícuota imaginaria).}

Cuando nos dan la concentración $c_{\mathrm{sol}}$ y no nos dan el volumen de 
la disolución, podemos asumir un volumen arbitrario que nos convenga para 
hacer los cálculos. Pero ¿por qué se puede hacer esa suposición? ¿De dónde 
sale ese volumen?

La respuesta es que estamos tomando una \textbf{alícuota imaginaria}: una 
porción mental de la misma disolución homogénea. No estamos creando disolución 
de la nada. La disolución ya tiene esa concentración en cualquier punto, y al 
elegir un volumen solo seleccionamos una porción de un sistema que ya era 
uniforme.

Imaginemos una jarra de agua de limón bien mezclada. Tiene una cierta 
concentración, no sabemos cuál, pero si tomamos un vaso y le damos un trago, 
¿sabe distinto que si tomamos directo de la jarra? No. La concentración es la 
misma en cualquier punto de la jarra. Esa es la propiedad intensiva de una 
disolución homogénea: no depende de cuánto tomes.

\medskip
Con eso en mente, podemos asumir un volumen que nos convenga (por ejemplo, 
$1\,\mathrm{L}$) y hacer los cálculos sobre esa alícuota imaginaria. El 
resultado no cambia, porque la concentración es la misma sin importar el 
volumen que elijamos.

\textit{Alícuota imaginaria vs.\ alícuota real.} En una alícuota real (como en 
el ejercicio 1, donde se toman $25.0\,\mathrm{mL}$ de $200.0\,\mathrm{mL}$) hay 
separación física y el volumen está limitado por la cantidad de disolución 
disponible. En la alícuota imaginaria no hay separación física ni límite: solo 
se elige un volumen para simplificar las cuentas. La diferencia es operativa, 
no conceptual: ambas conservan las propiedades intensivas de la disolución.
\end{notaimportante}

Aplicando esto al ejercicio:

\begin{align*}
    &V_{\mathrm{dil}} = 1.000\,\mathrm{L} = 1000\,\mathrm{mL} 
        \quad \text{(alícuota imaginaria)} \\[4pt]
    &n_{\mathrm{sol}} = c_{\mathrm{sol}} \cdot V_{\mathrm{dil}}
        = (2.25\,\mathrm{mol/L})(1.000\,\mathrm{L})
        = 2.25\,\mathrm{mol} \\[4pt]
    &MM_{\ce{C4H8O}} = 4(12.01) + 8(1.008) + 16.00 = 72.10\,\mathrm{g/mol} \\[4pt]
    &m_{\mathrm{sol}} = n_{\mathrm{sol}} \cdot MM_{\ce{C4H8O}}
        = (2.25\,\mathrm{mol})(72.10\,\mathrm{g/mol})
        = 162.2\,\mathrm{g} \\[4pt]
    &m_{\mathrm{dil}} = \rho_{\mathrm{dil}} \cdot V_{\mathrm{dil}}
        = (0.89\,\mathrm{g/mL})(1000\,\mathrm{mL})
        = 890\,\mathrm{g} \\[4pt]
    &m_{\mathrm{dis}} = m_{\mathrm{dil}} - m_{\mathrm{sol}}
        = 890\,\mathrm{g} - 162.2\,\mathrm{g}
        = 727.8\,\mathrm{g} = 0.7278\,\mathrm{kg} \\[4pt]
    &b_{\mathrm{sol}} = \frac{n_{\mathrm{sol}}}{m_{\mathrm{dis}}}
        = \frac{2.25\,\mathrm{mol}}{0.7278\,\mathrm{kg}}
        = 3.09\,\mathrm{mol/kg}
\end{align*}

\begin{notaimportante}
\textbf{Sobre las masas calculadas.} Las masas obtenidas 
($890\,\mathrm{g}$, $162.2\,\mathrm{g}$, $727.8\,\mathrm{g}$) corresponden a la 
alícuota imaginaria de $1\,\mathrm{L}$, no a la disolución original. No son 
masas reales, porque desconocemos el volumen total de la disolución. Son pasos 
intermedios para llegar a la molalidad, que sí es real e independiente del 
volumen elegido. La densidad del agua no se usa porque la masa de disolvente se 
obtiene por diferencia.
\end{notaimportante}
\end{exercise}

\textbf{La alícuota imaginaria como herramienta general.}

En el ejercicio anterior apareció la idea de asumir un volumen cuando falta el 
dato extensivo. Esto no es exclusivo de ese ejercicio: \textbf{todas las 
expresiones de concentración son propiedades intensivas} y permiten la misma 
herramienta. Al elegir una alícuota imaginaria, conservamos las propiedades 
intensivas de la disolución original y podemos calcular sin necesidad del dato 
faltante.

Esto aplica a cualquier expresión de concentración:
\begin{itemize}
    \item Molaridad ($c_{\mathrm{sol}}$)
    \item Molalidad ($b_{\mathrm{sol}}$)
    \item Normalidad ($N$)
    \item Fracción molar ($\chi_i$)
    \item Concentración en peso ($\rho_{\mathrm{sol}}$)
    \item Porcentaje en masa (\%\,m/m)
    \item Porcentaje masa/volumen (\%\,m/v)
    \item Porcentaje volumen/volumen (\%\,v/v)
    \item Partes por millón (ppm)
\end{itemize}

Todas son intensivas. En cualquier ejercicio donde falte un dato extensivo, se 
puede asumir una base de cálculo.

\textbf{¿Qué base elegir?} Depende del dato que te den. La regla es elegir la 
base que haga que el dato se lea directamente, sin conversiones extra:

\begin{center}
\begin{tabular}{|l|l|}
\hline
\textbf{Dato proporcionado} & \textbf{Base de cálculo natural} \\
\hline
Molaridad (mol/L) & $1.000\,\mathrm{L}$ de disolución \\
Normalidad (eq/L) & $1.000\,\mathrm{L}$ de disolución \\
Molalidad (mol/kg) & $1.000\,\mathrm{kg}$ de disolvente \\
\%\,m/m & $100\,\mathrm{g}$ de disolución \\
\%\,m/v & $100\,\mathrm{mL}$ de disolución \\
\%\,v/v & $100\,\mathrm{mL}$ de disolución \\
ppm (m/m) & $1\,\mathrm{kg}$ de disolución \\
ppm (m/v) & $1\,\mathrm{L}$ de disolución \\
Densidad (g/mL) & $1.000\,\mathrm{mL}$ de disolución \\
\hline
\end{tabular}
\end{center}

\textbf{Qué se conserva y qué no.} La alícuota imaginaria conserva las 
propiedades intensivas (concentración, densidad, composición) pero no las 
extensivas (masa total, volumen total, mol totales). Las cantidades 
extensivas que calcules a partir de ella son hipotéticas: sirven como pasos 
intermedios, no representan la disolución real.

\textit{Resumen:} la alícuota imaginaria es una herramienta para simplificar 
cálculos cuando falta un dato extensivo. Se elige el volumen o masa que más 
convenga según el dato dado. El resultado final (una propiedad intensiva) es 
real e independiente de la base elegida. Las cantidades intermedias 
extensivas son ficticias.

\begin{exercise}{7}{}
¿Cuál es la normalidad del \ce{H2SO4} cuya etiqueta dice $18\%\,m/m$ y 
densidad $1.15\,\mathrm{g/mL}$?

Solución:

\begin{align*}
    &N = ? \\[4pt]
    &\%\,m/m = 18\% \\[4pt]
    &\rho_{\mathrm{dil}} = 1.15\,\mathrm{g/mL} \\[4pt]
    &MM_{\ce{H2SO4}} = 2(1.008) + 32.06 + 4(16.00) = 98.08\,\mathrm{g/mol} \\[4pt]
\end{align*}

Como ya vimos con las propiedades intensivas, y siendo cualquier expresión de 
concentración una de ellas, no depende de la cantidad de materia. Igual que en 
el ejercicio anterior, podemos asumir una cantidad que nos convenga. Para este 
ejercicio tomamos una \textbf{alícuota imaginaria} de $100\,\mathrm{g}$ de 
disolución, de los cuales $18\,\mathrm{g}$ son de soluto. Eso sale directamente 
del $18\%\,m/m$.

\begin{align*}
    &m_{\mathrm{dil}} = 100\,\mathrm{g} \\[4pt]
    &m_{\mathrm{sol}} = 18\,\mathrm{g} \\[4pt]
    &V_{\mathrm{dil}} = \frac{m_{\mathrm{dil}}}{\rho_{\mathrm{dil}}}
    = \frac{100\,\mathrm{g}}{1.15\,\mathrm{g/mL}}
    = 86.96\,\mathrm{mL} = 0.08696\,\mathrm{L}\\[4pt]
    &n_{\mathrm{sol}} = \frac{m_{\mathrm{sol}}}{MM_{\mathrm{sol}}}
    = \frac{18\,\mathrm{g} }{98.08\,\mathrm{g/mol}}
    = 0.1836\,\mathrm{mol}
\end{align*}

\begin{center}
\schemestart
\ce{H2SO4}
\arrow{->}
\ce{2H^{\color{catred}+} + SO4^{\color{anblue}2-}}
\schemestop
\end{center}

\begin{align*}
    z_{\ce{H2SO4}} &= 2\,\mathrm{eq/mol} \\[4pt]
    n_{eq} &= z_{\ce{H2SO4}} \cdot n_{\mathrm{sol}}
    = 2\,\mathrm{eq/mol} \cdot 0.1836\,\mathrm{mol}
    = 0.3672\,\mathrm{eq}\\[4pt]
    N &= \frac{n_{eq}}{V_{\mathrm{dil}}}
    = \frac{0.3672\,\mathrm{eq}}{0.08696\,\mathrm{L}}
    =4.22\,\mathrm{eq/L}
\end{align*}
\end{exercise}

\begin{exercise}{8}{}
Una disolución contiene $125\,\mathrm{g}$ de fenol (\ce{C6H6O}) en 
$0.10\,\mathrm{kg}$ de agua. ¿Cuál es la fracción molar del fenol? 
¿Y del agua?

Solución:
\begin{align*}
    &\chi_{\ce{C6H6O}} = ? \\[4pt]
    &\chi_{\ce{H2O}} = ? \\[4pt]
    &m_{\ce{C6H6O}} = 125\,\mathrm{g} \\[4pt]
    &m_{\ce{H2O}} = 0.10\,\mathrm{kg} = 100\,\mathrm{g} \\[4pt]
    &MM_{\ce{C6H6O}} = 6(12.01) + 6(1.008) + 16.00 = 94.11\,\mathrm{g/mol} \\[4pt]
    &MM_{\ce{H2O}} = 2(1.008) + 16.00 = 18.016\,\mathrm{g/mol} \\[4pt]
    &n_{\ce{C6H6O}} = \frac{m_{\ce{C6H6O}}}{MM_{\ce{C6H6O}}}
    = \frac{125\,\mathrm{g}}{94.11\,\mathrm{g/mol}}
    = 1.33\,\mathrm{mol} \\[4pt]
    &n_{\ce{H2O}} = \frac{m_{\ce{H2O}}}{MM_{\ce{H2O}}}
    = \frac{100\,\mathrm{g}}{18.016\,\mathrm{g/mol}}
    = 5.55\,\mathrm{mol} \\[4pt]
    &n_{total} = n_{\ce{C6H6O}} + n_{\ce{H2O}}
    = 1.33\,\mathrm{mol} + 5.55\,\mathrm{mol}
    = 6.88\,\mathrm{mol} \\[4pt]
    &\chi_{\ce{C6H6O}} = \frac{n_{\ce{C6H6O}}}{n_{total}}
    = \frac{1.33\,\mathrm{mol}}{6.88\,\mathrm{mol}}
    = 0.193 \\[4pt]
    &\chi_{\ce{H2O}} = 1 - \chi_{\ce{C6H6O}}
    = 1 - 0.193 
    = 0.807
\end{align*}
\end{exercise}

\phantomsection
\addcontentsline{toc}{subsubsection}{Clase 4 -- 27/08/2026}
\subsubsection*{\textcolor{blueunam}{Clase 4 -- 27/08/2026: Ejercicios integradores de concentración}}

\begin{exercise}{9}{}
El Reglamento de Tránsito prohíbe conducir vehículos cuando el conductor 
presenta una cantidad de alcohol puro en sangre superior a $0.8\,\mathrm{g/L}$, 
o su equivalente en aire exhalado $0.4\,\mathrm{mg/L}$. Determina si un hombre 
de $75\,\mathrm{kg}$ puede conducir después de haber ingerido 4 cervezas con 
graduación alcohólica de $5\%\,v/v$. Toma en cuenta que, en promedio, el volumen 
de sangre, en L, de un adulto corresponde a un doceavo de su masa corporal, y 
que se estima que $15\%$ del alcohol ingerido llega a la sangre. Considera que 
la densidad del etanol es $0.789\,\mathrm{g/mL}$ y que una cerveza tiene un 
volumen de $355\,\mathrm{mL}$.

Solución:

\begin{notaimportante}
El enunciado dice ``alcohol'' genéricamente, pero la densidad dada 
($0.789\,\mathrm{g/mL}$) y la graduación de bebidas corresponden al etanol 
(\ce{C2H5OH}). En estos cálculos, ``alcohol'' = etanol.
\end{notaimportante}

\begin{align*}
    &c_{\text{etanol en sangre}} = ? \\[4pt]
    &c_{\text{límite}} = 0.8\,\mathrm{g/L} \\[4pt]
    &m_{\text{hombre}} = 75\,\mathrm{kg} \\[4pt]
    &\#_{\text{cervezas}} = 4 \\[4pt]
    &\%\,v/v = 5\% \\[4pt]
    &V_{\text{cerveza}} = 355\,\mathrm{mL} \\[4pt]
    &\rho_{\text{etanol}} = 0.789\,\mathrm{g/mL} \\[4pt]
    &\text{Porcentaje a sangre} = 15\% \\[4pt]
    &V_{\text{sangre}} = \frac{m_{\text{hombre}}}{12} 
    = \frac{75\,\mathrm{kg}}{12}
    = 6.25\,\mathrm{L}
\end{align*}

Ejercicios como estos pueden parecer complicados por la gran cantidad de datos, 
pero la clave es organizarlos, ver qué tenemos y ver qué queremos. En este caso, 
queremos la concentración de etanol en sangre, es decir, en unidades: la masa 
de soluto (etanol) que llega a la sangre en gramos y el volumen de disolución 
(sangre) en litros. El objetivo es calcular esa masa y dividirla entre el volumen 
de sangre. Parece obvio, pero verlo así ayuda mucho porque nos permite ver qué 
tenemos que hacer: una serie de conversiones para llegar a esas unidades. \\

Para el primer paso necesitamos saber cuánto etanol puro hay en las 4 cervezas. 
El dato es la graduación alcohólica, $5\%\,v/v$. Por definición, el $\%\,v/v$ 
es el volumen de soluto por cada 100 volúmenes de disolución. Entonces $5\%\,v/v$ 
significa: 5 mL de etanol por cada 100 mL de cerveza.

Esa relación es una concentración, y como toda concentración, es una propiedad 
intensiva. Imagina una alícuota imaginaria de exactamente $100\,\mathrm{mL}$ de 
cerveza: en ella habría $5\,\mathrm{mL}$ de etanol. Y como la concentración no 
depende del volumen, esa proporción se conserva en cualquier otra cantidad de 
cerveza. Si tomas $355\,\mathrm{mL}$ (una cerveza), hay 
$355 \times 5/100 = 17.75\,\mathrm{mL}$ de etanol. Si tomas $1420\,\mathrm{mL}$ 
(las 4 cervezas), hay $1420 \times 5/100 = 71\,\mathrm{mL}$ de etanol.

Es la misma idea de la alícuota: la proporción se conserva en cualquier volumen. 
Por eso podemos multiplicar directamente el volumen total por la fracción. A esa 
fracción la llamamos \(f\):

\[
f = \frac{5\,\mathrm{mL}_{\text{etanol}}}{100\,\mathrm{mL}_{\text{cerveza}}}
\]

Y así el cálculo queda:

\begin{align*}
    V_{\text{cerveza total}} &= \#_{\text{cervezas}} \cdot V_{\text{cerveza}}
    = 4 \cdot 355\,\mathrm{mL} 
    = 1420\,\mathrm{mL} = 1.420\,\mathrm{L} \\[4pt]
    V_{\text{etanol puro}} &= V_{\text{cerveza total}} \cdot f \\[4pt]
    &= 1420\,\mathrm{mL}_{\text{cerveza}} \cdot 
    \frac{5\,\mathrm{mL}_{\text{etanol}}}{100\,\mathrm{mL}_{\text{cerveza}}} \\[4pt]
    &= 71\,\mathrm{mL}_{\text{etanol}}
\end{align*}

\begin{align*}
    m_{\text{etanol}} &= V_{\text{etanol puro}} \cdot \rho_{\text{etanol}}
    = 71\,\mathrm{mL} \cdot 0.789\,\mathrm{g/mL}
    = 56.019\,\mathrm{g} \\[4pt]
    m_{\text{etanol en sangre}} &= m_{\text{etanol}} \cdot 
    \text{Porcentaje a sangre}
    = 56.019\,\mathrm{g} \cdot \frac{15}{100}
    = 8.402\,\mathrm{g} \\[4pt]
    c_{\text{etanol en sangre}} &= \frac{m_{\text{etanol en sangre}}}{V_{\text{sangre}}}
    = \frac{8.402\,\mathrm{g}}{6.25\,\mathrm{L}}
    = 1.344\,\mathrm{g/L}
\end{align*}

\textbf{Conclusión:} $c_{\text{etanol en sangre}} = 1.344\,\mathrm{g/L}$, que 
supera el límite legal de $0.8\,\mathrm{g/L}$. El hombre \textbf{no puede 
conducir}.
\end{exercise}

\phantomsection
\addcontentsline{toc}{subsubsection}{Clase 5 -- 01/09/2026}
\subsubsection*{\textcolor{blueunam}{Clase 5 -- 01/09/2026: Reacciones químicas y estequiometría}}

\textbf{Reacción química:} proceso en el que los átomos de las sustancias 
iniciales se reorganizan para formar sustancias nuevas. Esto implica romper 
enlaces en los reactivos y formar enlaces en los productos.

\begin{center}
\schemestart
$\underbrace{\ce{aA} + \ce{bB}}_{\text{Reactivos}}$
\arrow{->}
$\underbrace{\ce{cC} + \ce{dD}}_{\text{Productos}}$
\schemestop
\end{center}

Donde a, b, c y d son los coeficientes estequiométricos, que indican la proporción
en la que los reactivos se combinan y los productos se forman.

\begin{center}
\schemestart
$\underbrace{\ce{N2} + \ce{3H2}}_{\text{Reactivos}}$
\arrow{->}
$\underbrace{\ce{2NH3}}_{\text{Productos}}$
\schemestop
\end{center}

Las reacciones tienen la proporción de los reactivos y productos determinada 
por los coeficientes estequiométricos y esta proporción se mantiene sin importar 
la cantidad de materia que se use, es decir, podemos hablar tanto de moléculas individuales 
como de mol de cada especie. \\

\textbf{Nivel microscópico:}
\begin{center}
\schemestart
\chemfig{N~N}
\+
\chemfig{H-H}
\+
\chemfig{H-H}
\+
\chemfig{H-H}
\arrow{->}
\chemfig{N(-[:90]H)(-[:210]H)(-[:330]H)}
\+
\chemfig{N(-[:90]H)(-[:210]H)(-[:330]H)}
\schemestop
\end{center}

\textbf{Nivel macroscópico:}
\begin{center}
\schemestart
\chemname{\chemfig{N~N}}{$1\,\mathrm{mol}$}
\+
\chemname{3\,\chemfig{H-H}}{$3\,\mathrm{mol}$}
\arrow{->}
\chemname{2\,\chemfig{N(-[:90]H)(-[:210]H)(-[:330]H)}}{$2\,\mathrm{mol}$}
\schemestop
\end{center}

\textbf{Balance de ecuaciones químicas:}
\begin{enumerate}
    \item \textbf{Por tanteo:} ajustando los coeficientes estequiométricos 
    hasta que el número de átomos de cada elemento sea el mismo en ambos lados 
    de la ecuación.
    \begin{center}
    \schemestart
    \ce{2NaI}
    \+
    \ce{Br2}
    \arrow{->}
    \ce{2NaBr}
    \+
    \ce{I2}
    \schemestop
    \end{center}

    \item \textbf{Por el método algebraico:} escribiendo una ecuación para cada 
    elemento y resolviendo el sistema de ecuaciones resultante.
    \begin{enumerate}[label=Caso \arabic*]
        \item \(n\) incógnitas y \(n\) ecuaciones: sistema determinado, una 
        única solución.
        \item \(n\) incógnitas y \(n-1\) ecuaciones: sistema indeterminado, 
        infinitas soluciones. Se puede fijar un coeficiente arbitrario y resolver 
        el resto.
        \item \(n\) incógnitas y \(n+1\) o más ecuaciones: sistema 
        sobredeterminado. En general no tiene solución única; si la tiene, es 
        porque las ecuaciones extra son dependientes de las anteriores.
    \end{enumerate}

    \item \textbf{Método ion-electrón:} se balancea por separado la oxidación 
    y la reducción como semirreacciones, y luego se combinan.
\end{enumerate}

\textbf{Método por tanteo.}

El tanteo no es adivinar. Es seguir un orden que minimiza los reajustes. 
La estrategia:
\begin{enumerate}
    \item Empezar por el compuesto más complejo (el que tiene más elementos).
    \item Balancear primero metales, luego no metales.
    \item Dejar H y O para el final.
    \item Si aparece una fracción, multiplicar toda la ecuación por el 
    denominador.
\end{enumerate}

\textbf{Ejemplo:} balancear la combustión del octano.

\begin{center}
\schemestart
\ce{C8H18}
\+
\ce{O2}
\arrow{->}
\ce{CO2}
\+
\ce{H2O}
\schemestop
\end{center}

\textit{Paso 1.} El compuesto más complejo es \ce{C8H18}. Empezamos por él.
Ya tiene coeficiente 1 implícito.

\textit{Paso 2.} Balancear C. Hay 8 C a la izquierda, así que necesitamos 8 
\ce{CO2} a la derecha:

\begin{center}
\schemestart
\ce{C8H18}
\+
\ce{O2}
\arrow{->}
8\,\ce{CO2}
\+
\ce{H2O}
\schemestop
\end{center}

\textit{Paso 3.} Balancear H. Hay 18 H a la izquierda, así que necesitamos 
9 \ce{H2O} a la derecha (cada \ce{H2O} tiene 2 H, $18/2 = 9$):

\begin{center}
\schemestart
\ce{C8H18}
\+
\ce{O2}
\arrow{->}
8\,\ce{CO2}
\+
9\,\ce{H2O}
\schemestop
\end{center}

\textit{Paso 4.} Balancear O al final. A la derecha hay:
$8 \times 2 = 16$ O del \ce{CO2}, más $9 \times 1 = 9$ O del \ce{H2O}. 
Total: 25 O. A la izquierda, cada \ce{O2} aporta 2 O, así que necesitamos 
$25/2 = 12.5$ \ce{O2}:

\begin{center}
\schemestart
\ce{C8H18}
\+
12.5\,\ce{O2}
\arrow{->}
8\,\ce{CO2}
\+
9\,\ce{H2O}
\schemestop
\end{center}

\textit{Paso 5.} El coeficiente $12.5$ es fraccionario. Multiplicamos toda 
la ecuación por 2 para tener enteros:

\begin{center}
\schemestart
2\,\ce{C8H18}
\+
25\,\ce{O2}
\arrow{->}
16\,\ce{CO2}
\+
18\,\ce{H2O}
\schemestop
\end{center}

\textbf{Verificación:} 
C: $2 \times 8 = 16$ a la izquierda, $16$ a la derecha. 
H: $2 \times 18 = 36$ a la izquierda, $18 \times 2 = 36$ a la derecha. 
O: $25 \times 2 = 50$ a la izquierda, $16 \times 2 + 18 \times 1 = 50$ a la 
derecha. Balanceada.